% Chapter 1, typeset from lectures/L01: the README is the opener and the closing review, and each
% appendix is a section, in the same order and with the same subsections.

\chapter{The Register Package and the Peripheral Top}
\label{c1:ch}

{\large\itshape Top-down: the shell the whole peripheral fills in.\par}

\begin{chapterbox}{In this chapter}
\begin{itemize}
  \item \textbf{The provided register map, \code{uart_def}}: the seven register indices and the
    \code{STATUS}, \code{CTRL} and \code{ERROR_FLAGS} bit positions from Part~2 of the protocol
    (\secref{proto:part2}), mirrored on the C++ side in \file{register_map.hpp}.
  \item \textbf{\code{uart_top} itself}: the entity as a positional port contract, the provided
    \code{reset_sync} that asserts asynchronously and releases synchronously, and the provided SPI
    transport behind it.
  \item \textbf{Declaring signals as you need them}: each internal signal comes from reading the
    entity of the block it connects to, so the top grows one block's worth of wiring per chapter
    rather than starting with a list of names for modules that do not exist yet.
  \item \textbf{What the shape costs and buys}: what a positional contract cannot catch, and how the
    top grows across Chapters~2 to~5 until the system testbench can finally run.
\end{itemize}
\end{chapterbox}

This chapter starts at the top. Before a single datapath block exists we build the peripheral's
shell, because the shell is what fixes every block's interface: once \code{uart_top} says how it
instantiates \code{baud_gen}, the ports \code{baud_gen} must expose are settled, and the rest of
Part~I is filling that shell in. After the chapter you should be able to read the provided
register-map package and explain why its bit-\textbf{position} constants must match the protocol;
build the \code{uart_top} skeleton so that it analyzes cleanly with no datapath blocks present; and
explain why fixing the top first fixes each block's port contract, and why the system testbench is
skipped rather than failed until \chapref{5}.

\textbf{Before you start}, you need digital design with VHDL: \code{entity} and
\code{architecture}, structural instantiation, clocked processes, and analyzing a design with GHDL.
Read the register map (\secref{proto:part2}) and the SPI transport (\secref{proto:part3}) in the
protocol appendix; you instantiate the transport, but you do not write it.

Read \secref{c1:sec:uart_def} first, for the register map you will wire against, then
\secref{c1:sec:uart_top}, which is the reasoning behind \code{uart_top} and the design you build
from. The exercises in \secref{c1:sec:exercises} carry the build itself, so skim them as well and
know where their tables are. Do open \file{reset_sync.vhd}, \file{spi_slave.vhd} and
\file{spi_reg_bridge.vhd} in \file{hw/}: their entities are where the signals you have to declare
come from, and reading a port list to work out what wiring it implies is the habit this chapter is
teaching. The exercises tabulate the same ports, but as a check on your reading rather than a
substitute for it. Finally, read the system testbench \file{uart_top_tb.vhd} to see the contract the
top must eventually satisfy; you will not run it until \chapref{5}.

By the end of the chapter \code{uart_top} analyzes, though there is no testbench to run yet:

\begin{shell}
cd hw
ghdl -a --std=93 uart_def.vhd reset_sync.vhd spi_slave.vhd spi_reg_bridge.vhd uart_top.vhd
\end{shell}

\noindent Then, from the repository root, run \code{make build-vhdl} and note that
\code{uart_top_tb} reports \textbf{skipped}: its datapath blocks do not exist yet.

In the exercises, Exercise~\ref{c1:ex:2} is the build, decomposed into five steps: the entity, the
signals this chapter actually wires, the provided \code{reset_sync}, the two transport blocks, and
the check. Exercise~\ref{c1:ex:1} is the reading: the register map, positions rather than masks, and
when a swapped index would surface. Exercise~\ref{c1:ex:3} is the reasoning, about positional
binding, the reset synchronizer, and a register bus that waits four chapters for its bank.

% Section 1.1, from lectures/L01/appendix/a_uart_def.md, with the package itself typeset straight
% from hw/uart_def.vhd.

\section{\code{uart_def.vhd} (provided)}
\label{c1:sec:uart_def}\label{c1:app:a}
\code{uart_def} is the register map in VHDL form: one package of named constants for the seven
register indices and the meaning of each \code{STATUS}, \code{CTRL} and \code{ERROR_FLAGS} bit,
from Part~2 of the protocol (\secref{proto:part2}). Like the SPI transport, \textbf{it is given to
you}; you write every module in this course, but not the map. It exists so that the register bank
(\code{uart_regs}, \chapref{5}) and the testbenches can name registers and bits instead of writing
bare numbers, and so that the two sides of the wire cannot drift apart.

The C++ driver carries the same map in \file{register_map.hpp}, transcribed once on each side of the
wire. Both store \textbf{bit positions}, not masks, so \code{STATUS} bit~1 is the constant \code{1}
in both languages. The C++ side declares only the subset the driver actually uses, under the
\code{reg::}, \code{status::}, \code{ctrl::} and \code{error::} namespaces; where a name exists on
both sides, the value is identical. You do not edit \code{uart_def}, but read it closely: a module
that indexes the wrong bit is a bug you would only find on the bench.

\subsection{The constants}
\label{c1:sec:constants}
Everything is a \code{natural} in \code{package uart_def}. The register indices are the offset
divided by 4; the bit constants are positions into a 32-bit register word.

\begin{booktable}{@{}p{7.4em}p{8.6em}rL@{}}
  \thead{Group} & \thead{Constant} & \thead{Value} & \thead{Meaning} \\
  \midrule
  Register index & \code{REG_STATUS} & 0 & Status register. \\
   & \code{REG_CTRL} & 1 & Control register. \\
   & \code{REG_BAUD_DIV} & 2 & Baud divider. \\
   & \code{REG_TX_DATA} & 3 & Push a byte into the TX FIFO. \\
   & \code{REG_RX_DATA} & 4 & Front byte of the RX FIFO (a pure read). \\
   & \code{REG_RX_POP} & 5 & Advance the RX FIFO. \\
   & \code{REG_ERR_FLAGS} & 6 & Error flag register. \\
  \midrule
  \code{STATUS} bit & \code{ST_TX_READY} & 0 & The TX FIFO is not full. \\
   & \code{ST_RX_VALID} & 1 & The RX FIFO is not empty. \\
   & \code{ST_ERROR} & 2 & One or more error flags are set. \\
   & \code{ST_TX_IDLE} & 3 & The TX FIFO is empty and the line is idle. \\
  \midrule
  \code{CTRL} bit & \code{CT_ENABLE} & 0 & Enable the peripheral. \\
   & \code{CT_PARITY_LO} & 1 & Parity select, low bit. \\
   & \code{CT_PARITY_HI} & 2 & Parity select, high bit (00 none, 01 even, 10 odd). \\
   & \code{CT_STOP} & 3 & 0 for 1 stop bit, 1 for 2 stop bits. \\
   & \code{CT_RX_IRQ} & 4 & RX-valid IRQ mask. \\
   & \code{CT_TX_IRQ} & 5 & TX-ready IRQ mask. \\
  \midrule
  \code{ERROR_FLAGS} bit & \code{ER_FRAMING} & 0 & Framing error. \\
   & \code{ER_PARITY} & 1 & Parity error. \\
   & \code{ER_OVERRUN} & 2 & Overrun. \\
\end{booktable}

Because they are positions, you index a \code{std_logic_vector} with them directly, for example
\code{status(ST_RX_VALID)} for the RX-valid bit, exactly as \code{1U << status::RX_VALID} forms the
same mask on the C++ side. Laid out as words, the three flag registers are shown in
Figure~\ref{c1:fig:reg_fields}.

\bookfigure{reg_fields}{The bit fields of \code{STATUS}, \code{CTRL} and \code{ERROR_FLAGS}}
  {c1:fig:reg_fields}

\code{CTRL} is the only one with bits this course does not act on, and \code{ERROR_FLAGS} the only
one where two of the three have no producer yet; both are noted where they are built, in
\secref{c5:app:a}.

\subsection{The \code{to_hex} helper}
\label{c1:sec:to_hex}
The package also carries one small convenience for the testbenches: \code{to_hex}, which renders a
\code{std_logic_vector} as an uppercase, zero-padded hex string. It pads to whole nibbles taken from
the argument's own length, so \code{x"A5"} becomes \code{"A5"} and a 12-bit bus becomes three
digits. It lets a testbench name the offending value in a failure message,

\begin{vhdlcode}
report "rx_data byte 1 wrong: expected 0xA5, got 0x" & to_hex(captured_rx) & "!"
    severity error;
\end{vhdlcode}

\noindent rather than a bare ``mismatch''. It is the one piece of behaviour in an otherwise pure
constants package: declared in the package, defined in the package body, and called only by
testbenches. No synthesizable module uses it.

The whole package, constants and helper, is short enough to print as it stands in the repository:

\vhdlfile{hw/uart_def.vhd}

\subsection{Where it fits}
\label{c1:sec:uart_def_fits}
\code{uart_def} sits underneath everything that speaks in registers. \code{uart_regs}
(\chapref{5}) is the module that actually reads and writes these bits; the package only names them.
\code{uart_top} never uses it at all, because the top wires the register bus without interpreting
anything carried on it. Six testbenches analyze against the package, most of them only for
\code{to_hex}.

On the far side of the SPI wire, the driver you build from \chapref{6} on reaches the same
registers through the same positions. That is the point of transcribing the map twice rather than
inventing it twice: the protocol is the single source, and each language is checked against it
rather than against the other.

\subsection{What's ahead}
\secref{c1:sec:uart_top} is the peripheral top, \code{uart_top}, the structural skeleton built
against this package. The exercises in \secref{c1:sec:exercises} read \code{uart_def}, then build
\code{uart_top} from the ground up. The datapath and register modules the top composes arrive in
Chapters~2 to~5, and are instantiated into \code{uart_top} as each one lands.

% Section 1.2, from lectures/L01/appendix/b_uart_top.md, with the provided reset synchronizer
% typeset straight from hw/reset_sync.vhd.

\section{Designing \code{uart_top.vhd}}
\label{c1:sec:uart_top}\label{c1:app:b}
\code{uart_top} is the peripheral: the datapath (\code{baud_gen}, \code{uart_tx}, \code{uart_rx}),
the register bank (\code{uart_regs}), and the provided SPI transport (\code{spi_slave},
\code{spi_reg_bridge}) that lets an off-chip processor reach the registers. It is built
\textbf{first}, top-down, before any of those blocks exist. That is the point of starting here: the
top defines every block's interface, so when \code{baud_gen} arrives in \chapref{2} or
\code{uart_regs} in \chapref{5}, the ports each one must expose are already fixed by how
\code{uart_top} instantiates it.

What this chapter produces is the top's \emph{skeleton}: the entity, the provided
\code{reset_sync}, the provided transport, and the handful of signals those three need. It is
wiring, almost entirely: the pieces of real logic this file eventually holds, the TX feeder and one
type conversion, are written later, in the chapters that bring the signals they act on. The
datapath and register blocks do not exist yet; each is instantiated here as it is created across
Chapters~2 to~5, and the system testbench turns green once the last of them lands.

The signals follow the same rule. Rather than declaring names for modules that will not exist for
weeks, each chapter reads the entity of the block it is about to instantiate and declares exactly
the signals that block's ports imply. Two conventions keep it readable: a \textbf{\code{spi_}
prefix} for signals carrying SPI traffic between the two transport blocks, and an \textbf{\code{_s2}
suffix} for anything that has been through a two-flop synchronizer, which is why every submodule
expects \code{reset_s2_n}.

This section is the reasoning behind the file. The exercises in \secref{c1:sec:exercises} are where
it is actually typed, port by port and signal by signal.

\subsection{Interface}
\label{c1:sec:uart_top_interface}

\bookfigure{uart_top}{Module \code{uart_top}}{c1:fig:uart_top}

\begin{booktable}{@{}p{8.6em}p{2em}p{6em}L@{}}
  \thead{Port} & \thead{Dir} & \thead{Type} & \thead{Meaning} \\
  \midrule
  \code{clock} & in & \code{std_logic} & 50\,MHz system clock. \\
  \code{reset_n} & in & \code{std_logic} & Asynchronous reset, active low (a button or a
    supervisor). \\
  \code{sclk}, \code{mosi}, \code{ss} & in & \code{std_logic} & SPI from the processor. \\
  \code{rx} & in & \code{std_logic} & The UART serial input. \\
  \code{miso} & out & \code{std_logic} & SPI to the processor. \\
  \code{tx} & out & \code{std_logic} & The UART serial output. \\
\end{booktable}

The ports are declared all inputs first, then all outputs; \code{uart_top_tb} binds them
positionally, so the order above is the contract. Nothing in this course associates a port by
name, which means a transposed pair of same-type pins is invisible to the analyzer, and stays
invisible until the system testbench runs at the end of \chapref{5}.

\code{uart_top} is the only block that takes \code{reset_n}, the raw asynchronous reset, from
outside. The one block inside it that also takes \code{reset_n} is \code{reset_sync}, because
manufacturing a clean reset is its whole job; every other block takes \code{reset_s2_n}, the
synchronized version it produces.

\subsection{Skeleton behaviour (wiring, and one placeholder)}
\label{c1:sec:uart_top_behaviour}
The transport comes first, because it is what makes a register bus exist at all. \code{spi_slave}
and \code{spi_reg_bridge} are given to you, testbench-verified; their wire protocol is Part~3 of the
protocol (\secref{proto:part3}), and their port lists are shown in Exercise~\ref{c1:ex:2}. You
instantiate them; you do not write them. \code{spi_slave} is the byte engine: it shifts bytes in and
out on \code{sclk} and reports each completed one. \code{spi_reg_bridge} assembles those bytes into
the five-byte register transactions the protocol defines, drives \code{reg_addr}, \code{reg_wdata}
and \code{reg_write}, and reads back whatever \code{reg_rdata} presents to it. The bank that will
answer on that bus arrives in \chapref{5}; until then the bus exists with nothing behind it.

The \textbf{reset synchronizer} is provided, as \code{reset_sync}, and you instantiate it.
\code{reset_n} comes from off-chip, so it can be released at any instant relative to \code{clock}. A
release landing near a clock edge could let different flip-flops leave reset on different cycles,
and a design that starts half-reset can sit in a state its own logic never expects.
\code{reset_sync} asserts asynchronously, so the moment \code{reset_n} goes low everything resets
with no clock needed, and releases synchronously, two flip-flops later, so the release is lined up
with the clock. Its two internal flops stay inside it, so the only signal it adds to your top is
\code{reset_s2_n}. The whole module is short enough to print:

\vhdlfile{hw/reset_sync.vhd}

\noindent That output then feeds every submodule, and the blocks you write in Chapters~2 to~5 all
take it the same way the provided ones do: \code{reset_s2_n} goes in the process's sensitivity
list, so the assertion still needs no clock edge, while the \emph{release} is already lined up
because \code{reset_sync} lined it up. That is the house style throughout \file{hw/}, and it is what
makes an asynchronous assertion safe in the first place: the dangerous edge, the release, was dealt
with once, here.

Note that this job is not \code{sync}'s (\chapref{3}), even though that module is also two flops in
series. \code{sync} takes \code{reset_s2_n} as an \emph{input}, so it consumes the very signal
\code{reset_sync} produces, and a block that needs a clean reset cannot be the block that
manufactures one.

The \textbf{TX feeder} is written in \chapref{5}, once both signals it needs exist. \code{uart_tx}
(\chapref{2}) sends a byte when its \code{start} is pulsed, and the TX FIFO inside \code{uart_regs}
(\chapref{5}) holds the bytes waiting to go. The feeder is what connects them: \code{tx_load} is
FIFO-not-empty AND transmitter-not-busy, and \code{tx_pop} is that same signal, so the transmitter
latches the front byte and the FIFO advances in the same cycle. The instant the transmitter goes
busy, \code{tx_load} drops, so exactly one byte moves per idle transmitter, never two. The RX side
needs no feeder at all: the receiver's \code{valid} and \code{data_out} become \code{uart_regs}'
\code{rx_push} and \code{rx_byte}, and the bank pushes its own FIFO.

The other piece is a \textbf{type conversion}, written in \chapref{2} alongside \code{baud_gen}.
\code{baud_gen} takes its divider as a \code{natural range 1 to 65535}, but \code{BAUD_DIV} leaves
the register bank as a \code{std_logic_vector(15 downto 0)}, so \code{uart_top} converts it. That
conversion is guarded, because until \code{uart_regs} exists in \chapref{5} nothing drives the
vector, and \code{to_integer} on an all-\code{'U'} value warns and substitutes a zero of its own
choosing, which is a value that port's range does not accept. Picking a legal value yourself costs
one conditional expression, and this is the only place in \code{uart_top} where a
\code{std_logic_vector} becomes an integer.

The one thing this chapter does write is a \textbf{placeholder}: \code{reg_rdata} tied to zeros, so
the bridge reads a defined value rather than \code{'U'} while the register bank is missing. It is
deleted in \chapref{5}, when \code{uart_regs} starts driving that vector for real.

\subsection{What the testbench pins down}
\label{c1:sec:uart_top_tb}
\code{uart_top_tb} is the system testbench, and it does not run until the last block lands in
\chapref{5}. It drives the peripheral entirely over SPI, as the processor would, and loops \code{tx}
back into \code{rx}. A single byte written over SPI is therefore transmitted, received through the
loopback, pushed into the RX FIFO, and read back over SPI: an end-to-end check that the reset
synchronizer, the transport, the register bank, the FIFOs, the feeder and both datapath halves all
agree. It also writes \code{BAUD_DIV} and reads it back, and checks \code{STATUS}, so the register
plumbing is exercised alongside the datapath.

Its one line of instantiation is also why the entity looks the way it does. The testbench binds
\code{port map(clock, reset_n, sclk, mosi, ss, rx, miso, tx)}, and that order, not the drawing in
Figure~\ref{c1:fig:uart_top}, is what the port table records.

\subsection{Where it fits}
\label{c1:sec:uart_top_fits}
\code{uart_top} faces two directions. Outward, its eight pins are the FPGA boundary: \code{sclk},
\code{mosi}, \code{ss} and \code{miso} reach the ATmega328P on the other side of the wire, and the
driver built in Chapters~6 to~9 talks to exactly these ports; \code{rx} and \code{tx} are the
serial line itself. In \chapref{9} a board wrapper sits above it and maps those pins onto DE0-CV
package pins, which is the only thing that wrapper does.

Inward, it is the one file revisited in every VHDL chapter, and each visit adds a block, the
signals that block's entity implies, and any glue those signals now make possible. This chapter
leaves it with the entity, \code{reset_sync}, the transport and a register bus nothing answers.
\chapref{2} adds \code{baud_gen} and \code{uart_tx}, and with them the \code{BAUD_DIV} conversion,
so the top can transmit. \chapref{4} adds \code{sync} on the \code{rx} pin, producing \code{rx_s2},
and \code{uart_rx} behind it, so the receiver never sees the raw asynchronous input. \chapref{5}
adds \code{uart_regs}, which owns both FIFOs, and the TX feeder that connects it to the transmitter;
with every block present, the system testbench elaborates for the first time. The rule that keeps
it compiling throughout is that you only ever instantiate an entity after you have created it.

\subsection{What's ahead}
The exercises in \secref{c1:sec:exercises} are, for \code{uart_top}, also the build sheet:
everything needed to type the file, from the entity down to the last internal signal, is there. In
\chapref{2} you build the first two datapath blocks, \code{baud_gen} and \code{uart_tx}, and
instantiate them here.

% Section 1.3, from lectures/L01/README.md: the objectives, the evaluation questions, and the next
% lecture.

\section{Review}
\label{c1:sec:review}

\needspace{6\baselineskip}
\subsection*{What you should be able to do now}
After this chapter, you should be able to:
\begin{itemize}
  \item \textbf{Read the provided register-map package} (\code{uart_def}) and explain why its
    \code{natural} bit-\textbf{position} constants must match the protocol, and why the C++
    \file{register_map.hpp} agrees with it value for value wherever a name exists on both sides.
  \item \textbf{Build the \code{uart_top} skeleton}, meaning the entity, the provided
    \code{reset_sync}, the provided transport, and the signals those three need, so that it analyzes
    cleanly with no datapath blocks present.
  \item \textbf{Derive the internal signals from the entities they connect}, rather than from a
    list, and apply the \code{spi_} prefix and \code{_s2} suffix conventions the provided modules
    already use.
  \item \textbf{Explain top-down structural design}: why fixing the top first fixes each block's
    port contract, and why the system testbench is gated, skipped rather than failed, until
    \chapref{5}.
\end{itemize}

\needspace{6\baselineskip}
\subsection*{Questions to test yourself}
You should be able to explain:
\begin{enumerate}
  \item Why build \code{uart_top} before any block it instantiates, and what does fixing the top
    first fix about every datapath module you write later?
  \item The reset is asserted asynchronously but released synchronously: what can go wrong if the
    release is not synchronized, and why can \code{sync} (\chapref{3}) not do this job?
  \item \code{uart_top_tb} binds \code{uart_top}'s ports positionally: what class of wiring mistake
    does that leave the analyzer unable to catch, and at which point in the book does it finally
    surface?
\end{enumerate}

\needspace{6\baselineskip}
\subsection*{Looking ahead}
The datapath begins in \chapref{2}: \code{baud_gen}, the shared time base, and \code{uart_tx}, the
transmitter, each verified against its own testbench and instantiated into the \code{uart_top} you
built here.

% Section 1.4, from lectures/L01/appendix/c_exercises.md.

\section{Exercises}
\label{c1:sec:exercises}\label{c1:app:c}

\begin{aside}
\textbf{How to work these.} Exercise~\ref{c1:ex:1} reinforces \secref{c1:sec:uart_def};
Exercises~\ref{c1:ex:2} and~\ref{c1:ex:3} go with \secref{c1:sec:uart_top}. \code{uart_def} is
\textbf{provided} and only needs reading. \code{uart_top} is yours, and Exercise~\ref{c1:ex:2} is
self-contained: everything needed to type \file{hw/uart_top.vhd}, from the entity down to the last
internal signal, is below. \secref{c1:sec:uart_top} is where the reasoning behind it lives, and
Exercise~\ref{c1:ex:3} asks for that reasoning back.

\textbf{How to check your work.} Neither module has a testbench of its own in this chapter.
\code{uart_def} is a package of constants, and \code{uart_top}'s system testbench,
\code{uart_top_tb}, stays skipped until its datapath exists in \chapref{5}. So the check here is
that \code{uart_top} \textbf{analyzes cleanly} against the transport it instantiates. Run GHDL from
\file{hw/}; the preface's \emph{Setting up} has the full flow and what \code{--assert-level=error}
does.

The answers to every part that is not code are in \secref{sol:c1}. Try each part before you read its
answer.
\end{aside}

% ------------------------------------------------------------------------------------------
\exercise{\code{uart_def}}{Reasoning}
\label{c1:ex:1}

\xpart{a} \code{uart_def.vhd} is provided. Read it against \secref{c1:sec:uart_def}, then, without
looking, name the seven register indices and the \code{STATUS}, \code{CTRL} and \code{ERROR_FLAGS}
bit positions. Analyze it, so that the names are in your \code{work} library for everything that
follows:

\begin{shell}
cd hw
ghdl -a --std=93 uart_def.vhd
\end{shell}

\noindent There is nothing to elaborate or run; a clean analysis means the package compiles and its
names are available to every module that later writes \code{use work.uart_def.all}.

\xpart{b} These are bit \textbf{positions}, not masks: \code{ST_RX_VALID} is \code{1}, not
\code{2}. Show how you would test the RX-valid bit of a status vector \code{status} using the
constant, and contrast it with how the C++ driver forms the same test
(\code{1U << status::RX_VALID}). Why does storing positions keep both sides identical to the
protocol's ``bit N'' wording?

\xpart{c} A teammate swaps two index values, setting \code{REG_RX_POP} to 6 and
\code{REG_ERR_FLAGS} to 5. Nothing in this chapter notices. At what point in the course would the
swap first surface as a failure, and in which testbench, if any? Say why it does not surface
before then.

% ------------------------------------------------------------------------------------------
\exercise{\code{uart_top}}{VHDL}
\label{c1:ex:2}
Build the skeleton in the order the file is written: the entity, the declarations, then the three
instantiations and the one placeholder assignment. Each part below carries what that step needs. Do
not instantiate \code{baud_gen}, \code{uart_tx}, \code{sync}, \code{uart_rx} or \code{uart_regs}
yet, and do not declare their signals: those blocks arrive in Chapters~2 to~5, and each brings its
own wiring with it.

\xpart{a}\parttitle{The entity.} Declare the eight ports in exactly this order. \code{uart_top_tb}
binds them by position, so the numbers are the contract; the names are yours.

\begin{narrowtable}{@{}rlll@{}}
  \thead{\#} & \thead{Port} & \thead{Dir} & \thead{Type} \\
  \midrule
  1 & \code{clock} & in & \code{std_logic} \\
  2 & \code{reset_n} & in & \code{std_logic} \\
  3 & \code{sclk} & in & \code{std_logic} \\
  4 & \code{mosi} & in & \code{std_logic} \\
  5 & \code{ss} & in & \code{std_logic} \\
  6 & \code{rx} & in & \code{std_logic} \\
  7 & \code{miso} & out & \code{std_logic} \\
  8 & \code{tx} & out & \code{std_logic} \\
\end{narrowtable}

The file needs \code{ieee.std_logic_1164} for \code{std_logic}. It does not need
\code{ieee.numeric_std} yet, since nothing here converts a vector to an integer, and it does
\textbf{not} need \code{use work.uart_def.all}: the top wires the register bus but never names a bit
on it, and \code{uart_regs} (\chapref{5}) is the only module that does. Follow the provided files'
house style: a banner comment listing the inputs and outputs, then
\code{architecture behaviour of uart_top is}.

\xpart{b}\parttitle{The internal signals.} Declare only the signals \textbf{this} chapter wires.
The method matters more than the list: open the entity of every module you are about to
instantiate, and for each port that does not connect straight to a \code{uart_top} port, declare a
signal of the same type to carry it. Here that means \code{reset_sync}, \code{spi_slave} and
\code{spi_reg_bridge}, and nothing else. Each later chapter adds its own signals the same way, when
it adds the block that needs them.

Two naming conventions run through the whole design, and both are worth adopting now. Signals that
carry SPI traffic between the two transport blocks take a \textbf{\code{spi_} prefix}. Signals that
have been through a two-flop synchronizer take an \textbf{\code{_s2} suffix}, which is why the
provided modules all expect \code{reset_s2_n} rather than \code{reset_n}.

\begin{booktable}{@{}p{6.8em}p{15.2em}p{7.6em}L@{}}
  \thead{Signal} & \thead{Type} & \thead{Driven by} & \thead{Read by} \\
  \midrule
  \code{reset_s2_n} & \code{std_logic} & \code{reset_sync}, in c) & every submodule \\
  \code{spi_rx_data} & \code{std_logic_vector(7 downto 0)} & \code{spi_slave} &
    \code{spi_reg_bridge} \\
  \code{spi_rx_valid} & \code{std_logic} & \code{spi_slave} & \code{spi_reg_bridge} \\
  \code{spi_ss_active} & \code{std_logic} & \code{spi_slave} & \code{spi_reg_bridge} \\
  \code{spi_tx_data} & \code{std_logic_vector(7 downto 0)} & \code{spi_reg_bridge} &
    \code{spi_slave} \\
  \code{reg_addr} & \code{std_logic_vector(3 downto 0)} & \code{spi_reg_bridge} &
    \code{uart_regs} (\chapref{5}) \\
  \code{reg_wdata} & \code{std_logic_vector(31 downto 0)} & \code{spi_reg_bridge} &
    \code{uart_regs} (\chapref{5}) \\
  \code{reg_write} & \code{std_logic} & \code{spi_reg_bridge} & \code{uart_regs} (\chapref{5}) \\
  \code{reg_rdata} & \code{std_logic_vector(31 downto 0)} & zeros, in d); then \code{uart_regs}
    (\chapref{5}) & \code{spi_reg_bridge} \\
\end{booktable}

The four register-bus signals are here even though the block that answers on them is four chapters
away, because \code{spi_reg_bridge} has ports for them and an instantiation must connect every port.
That is the register bus existing with nothing behind it, which is the shape this chapter is really
about.

\code{reg_addr} is four bits, not three, even though there are only seven registers: the protocol's
command byte reserves bits 3 to 0 for the index.

\xpart{c}\parttitle{The reset synchronizer.} \code{reset_n} arrives from off-chip, so it must
assert asynchronously and release synchronously, two flip-flops later. That job is done by the
provided \code{reset_sync}, which you instantiate rather than write; its two internal flops are its
own business, so the only signal it adds to your top is \code{reset_s2_n}. Everything inside
\code{uart_top} takes \code{reset_s2_n}, and nothing but this instance sees \code{reset_n} directly.

\bookfigure{reset_sync}{Module \code{reset_sync}}{c1:fig:reset_sync}

\begin{narrowtable}{@{}rlll@{}}
  \thead{\#} & \thead{\code{reset_sync} port} & \thead{Dir} & \thead{Connect to} \\
  \midrule
  1 & \code{clock} & in & \code{clock} \\
  2 & \code{reset_n} & in & \code{reset_n} \\
  3 & \code{reset_s2_n} & out & \code{reset_s2_n} \\
\end{narrowtable}

\begin{vhdlcode}
reset_sync: entity work.reset_sync
    port map(clock, reset_n, reset_s2_n);
\end{vhdlcode}

\xpart{d}\parttitle{The transport.} Instantiate \code{spi_slave} and \code{spi_reg_bridge}, both
positionally. The last column names a signal from part b) or a port from part a).

\bookfigure{spi_slave}{Module \code{spi_slave}}{c1:fig:spi_slave}

\begin{booktable}{@{}rp{7.4em}p{1.8em}p{15.2em}L@{}}
  \thead{\#} & \thead{\code{spi_slave} port} & \thead{Dir} & \thead{Type} & \thead{Connect to} \\
  \midrule
  1 & \code{clock} & in & \code{std_logic} & \code{clock} \\
  2 & \code{reset_s2_n} & in & \code{std_logic} & \code{reset_s2_n} \\
  3 & \code{sclk} & in & \code{std_logic} & \code{sclk} \\
  4 & \code{mosi} & in & \code{std_logic} & \code{mosi} \\
  5 & \code{ss} & in & \code{std_logic} & \code{ss} \\
  6 & \code{tx_data} & in & \code{std_logic_vector(7 downto 0)} & \code{spi_tx_data} \\
  7 & \code{miso} & out & \code{std_logic} & \code{miso} \\
  8 & \code{rx_data} & out & \code{std_logic_vector(7 downto 0)} & \code{spi_rx_data} \\
  9 & \code{rx_valid} & out & \code{std_logic} & \code{spi_rx_valid} \\
  10 & \code{ss_active} & out & \code{std_logic} & \code{spi_ss_active} \\
\end{booktable}

\bookfigure{spi_reg_bridge}{Module \code{spi_reg_bridge}}{c1:fig:spi_reg_bridge}

\begin{booktable}{@{}rp{9.8em}p{1.8em}p{15.2em}L@{}}
  \thead{\#} & \thead{\code{spi_reg_bridge} port} & \thead{Dir} & \thead{Type} &
    \thead{Connect to} \\
  \midrule
  1 & \code{clock} & in & \code{std_logic} & \code{clock} \\
  2 & \code{reset_s2_n} & in & \code{std_logic} & \code{reset_s2_n} \\
  3 & \code{ss_active} & in & \code{std_logic} & \code{spi_ss_active} \\
  4 & \code{rx_data} & in & \code{std_logic_vector(7 downto 0)} & \code{spi_rx_data} \\
  5 & \code{rx_valid} & in & \code{std_logic} & \code{spi_rx_valid} \\
  6 & \code{reg_rdata} & in & \code{std_logic_vector(31 downto 0)} & \code{reg_rdata} \\
  7 & \code{tx_data} & out & \code{std_logic_vector(7 downto 0)} & \code{spi_tx_data} \\
  8 & \code{reg_addr} & out & \code{std_logic_vector(3 downto 0)} & \code{reg_addr} \\
  9 & \code{reg_wdata} & out & \code{std_logic_vector(31 downto 0)} & \code{reg_wdata} \\
  10 & \code{reg_write} & out & \code{std_logic} & \code{reg_write} \\
\end{booktable}

The names that look shared are not one bus. \code{spi_slave}'s \code{rx_data} is an output and the
bridge's \code{rx_data} is an input, so \code{spi_rx_data} runs from the slave to the bridge;
\code{spi_tx_data} runs the other way, from the bridge's \code{tx_data} output into the slave's
\code{tx_data} input. Read them as two one-directional lanes and the wiring falls out.

Nothing drives \code{reg_rdata} yet, so give it a value:

\begin{vhdlcode}
reg_rdata <= (others => '0');
\end{vhdlcode}

\noindent The file would analyze without this; an undriven signal is legal. It is here so that the
bridge reads a defined value instead of \code{'U'} while the bank is missing. \textbf{Delete this
line in \chapref{5}}, when \code{uart_regs} starts driving \code{reg_rdata}. Leave it in and the
vector has two drivers: every \code{'1'} from the bank resolves against the \code{'0'} here to
\code{'X'}, and the system testbench fails on a read-back caused by a line you wrote four chapters
earlier.

\xpart{e}\parttitle{The check.} Analyze the skeleton against the package and the transport it
instantiates:

\begin{shell}
cd hw
ghdl -a --std=93 uart_def.vhd reset_sync.vhd spi_slave.vhd spi_reg_bridge.vhd uart_top.vhd
\end{shell}

\noindent Then run \code{make build-vhdl} from the repository root. Confirm that \code{uart_def} and
\code{uart_top} both analyze cleanly, that the two transport testbenches pass (they are the only
ones that can run this early, since they check provided modules), and that \code{uart_top_tb}
reports \textbf{skipped}. Which modules does the skip message name as missing, and which chapter
(2, 3, 4 or 5) builds each? One of the six is never instantiated by \code{uart_top} at all; find
it, and say why the build needs it anyway.

% ------------------------------------------------------------------------------------------
\exercise{What a clean analysis does not prove}{Reasoning}
\label{c1:ex:3}
\code{uart_top} analyzes, and almost nothing about it is proven. That gap is the lesson of this
chapter.

\xpart{a} Swap two same-type ports in your entity, say \code{sclk} and \code{mosi}, and re-analyze.
Does \code{ghdl -a} catch it? If not, at which point in the course does the mistake finally show up,
and what does that tell you about the cost of a positional contract?

\xpart{b} \code{reg_rdata} is tied to zeros and no register bank answers on the bus. What does an
SPI read transaction return with that placeholder in place, and why is that harmless until
\chapref{5}? What would it return instead if you left the placeholder out altogether?

\xpart{c} The provided \code{reset_sync} asserts asynchronously but releases synchronously. Suppose
it released asynchronously too, wiring \code{reset_n} straight to every block. Describe the failure
that a near-edge release could cause. Then explain why \code{sync} (\chapref{3}) cannot do this job,
given that it is also two flip-flops in series: look at its port list and say what it would need
that only \code{reset_sync} can give it.

\xpart{d} Why does building \code{uart_top} before any block it instantiates fix each block's port
list in advance? What has to be true about \code{baud_gen}'s ports in \chapref{2} for its
instantiation here to bind without you touching the top again?

